\section{Related Work}
\paragraph{Multimodal Representation Learning}
The fundamental solutions of current multimodal tasks focus on multimodal representation learning, which dedicates to extracting and integrating task-related information from the input signals of many modalities~\cite{atrey2010multimodal,ngiam2011multimodal}. 
Recently, multimodal fusion technique becomes the predominant method for expressive representation learning, which coalesces a set of multimodal inputs by 
% neural network 
mathematical operations (e.g., attention and Cartesian product)~\cite{liu2018efficient,tsai2019multimodal,mohla2020fusatnet,hazarika2020misa,han2021bi}.
Though showing exceptional performance on 
% many downstream multimodal 
those tasks, stacked attention architecture also consumes huge computational power and slows down the training and inference speed.
% Distinct from these works,
To alleviate this issue, we devise a fusion-free model for the MRHP task, which escapes from the conventional fusion-based routine.

\paragraph{Relation-based Learning}
The idea of relation-based learning was firstly applied in the few-shot image classification task~\cite{vanschoren2018meta,hospedales2020meta}.
\citet{vinyals2016matching}~employs quantified similarity values between the unseen test images and seen trained images to perform classification.
Prototypical networks~\cite{snell2017prototypical} and Relation Network~\citet{sung2018learning} further treats the correlation matrices between images and pre-computed prototypical feature vectors as logits and optimize them to improve the model's performance.
\citet{lifchitz2019dense}~substitutes the comparison target with implanting weights for better generalization ability.
Later achievements stemming from this theory encompass building up network structures for interactions between samples~\cite{garcia2017few,kim2019edge}, incorporating small-scale computation units like image pixels~\cite{chang2018pyramid,si2018dual,hou2019cross,min2021hypercorrelation} and adding correlation matrices as regularization terms~\cite{wertheimer2021few}.
In the multimodal scenario, semantic matching has also been chosen as the core task to pretrain large multimodal models~\cite{chen2020uniter,kim2021vilt,radford2021learning}.
We inherit this idea to develop a matching-based approach for the MRHP task. 
In our network, sorted matching scores between vectors of different scales and modalities are shaped into regression features. 
We will show this canonical formulation beats many fusion-based strong baselines. 

